\subsubsection{Phenomenological implications of neutrinos in extra dimensions --- arXiv:hep-ph/0107156}

\textbf{Authors:} Andr\'e de Gouv\^ea, Gian Francesco Giudice, Alessandro Strumia, Kazuhiro Tobe

\paragraph{主な主張}
複数の大きな余剰次元では軽いKKモードが混合全体を支配しないため、既知の太陽・大気ニュートリノ振動をKKモードで説明するよりも、重いKKモードによる間接効果の探索が有望である \cite{deGouvea:2001ExtraDimensions}。

\paragraph{新規性}
重いKKモードの積分消去から次元6の有効演算子を導き、ニュートリノ混合行列のnon-unitarityと荷電レプトンのフレーバー非保存過程を統一的に解析した。

\paragraph{位置付け}
多次元bulk neutrino模型におけるKK混合のUV依存性と低エネルギー有効理論を扱った基礎的な先行研究であり、低いKKモードによる追加的な振動現象の観測可能性を否定するものではない。
